\documentclass[review,12pt,authoryear]{elsarticle}
\input{jmacros.tex}
\journal{Agricultural Water Management}
\hypersetup{pdftitle={The reach of suction-triggered irrigation and the marginal fate of irrigation water},pdfauthor={Naziru Halilu}}
\newcommand{\hc}{h_{\mathrm{c}}}
\newcommand{\htr}{h_{\mathrm{t}}}
\begin{document}

\begin{frontmatter}
\title{The reach of suction-triggered irrigation and the marginal fate of irrigation water: closed forms from the van~Genuchten--Mualem parameters}

\author[aua,upna]{Naziru Halilu\corref{cor1}}
\ead{halilunaziru73@gmail.com}
\cortext[cor1]{Corresponding author. ORCID: 0009-0006-2324-9599.}
\affiliation[aua]{organization={Agricultural University of Athens},city={Athens},country={Greece}}
\affiliation[upna]{organization={Public University of Navarre},city={Pamplona},country={Spain}}

\begin{abstract}
Suction triggers decide when most sensor-scheduled fields are irrigated, and the van~Genuchten--Mualem functions describe how a soil holds and conducts water at any suction. Between the two lies an open question: which part of the retention curve can a crop actually reach at the sensor, so that a trigger set beyond it never fires? A second question is where each applied millimetre of water goes as the soil, the dose, the method and the season change. Both were answered from the dry-end asymptote of the Mualem conductivity and an open ensemble of Daisy simulations of ten soils over two decades. Equating the relative conductivity at the reach to one constant shared by all soils gave a closed-form reach that depends only on the van~Genuchten shape parameters and the pore-connectivity exponent. It predicted the reach of soils left out of the fit within about a factor of two, from sands to silty clay loams. Commonly tabulated triggers lay beyond the reach of sandy soils: they never fired and the crop stayed at its rainfed yield. Scaling each trigger by the reach placed the irrigation responses of all soils on one axis. Less than half of each irrigation millimetre was transpired; most of the rest evaporated and a smaller share drained. The drained share rose with the dose towards a soil-specific plateau on a dose scale common to all soils and close to one day of crop water use, and overhead irrigation added no transpiration.
\end{abstract}

\begin{highlights}
\item A closed form gives the reach of a suction trigger from van Genuchten parameters
\item It predicted the reach of held-out soils within a factor of 2.3
\item The shape of K(h), not its saturated level, sets where a trigger stops firing
\item Of each irrigation millimetre, 38\% was transpired, 44\% evaporated, 16\% drained
\item The drained share rose with dose on a soil-independent scale of about 4 mm
\end{highlights}

\begin{keyword}
irrigation trigger \sep matric-potential sensor \sep van~Genuchten--Mualem \sep deep percolation \sep irrigation dose \sep Daisy model
\end{keyword}
\end{frontmatter}

\linenumbers

\section{Introduction}
Since \citet{vanGenuchten1980}, the retention curve of a soil and its unsaturated conductivity have been written with three shape parameters, $\alpha$, $n$ and the Mualem connectivity $\ell$, and that closed form now sits inside almost every model of the unsaturated zone \citep{vanGenuchten1985,Nielsen1986,Schaap2001,Simunek2016,Vereecken2016,Vereecken2019}. On many fields irrigation is scheduled with sensors and a suction threshold: when a tensiometer or matric-potential sensor at a chosen depth reads more than a trigger $\htr$, water is applied \citep{Allen1998}. The triggers are taken from tables (commonly 300 to 600\,cm for field crops), and the retention curve is used, if at all, to translate them into water contents.

Two questions lie between these lines of work. The first concerns the curve itself. The van~Genuchten curve is defined for every suction, but a crop and a sensor visit only part of it. The wet end of that part has been given intrinsic definitions from the hydraulic parameters: field capacity as the state where internal drainage becomes negligible \citep{Twarakavi2009,Assouline2014}. The dry end has no equivalent closed form for control: which suction can roots actually impose at the depth of the sensor, so that a trigger set beyond it is never reached? \citet{Gardner1960} showed that uptake becomes limited by the conductance of the soil around the roots long before the soil reaches the wilting point, a limit that root-uptake models express as a reduction of uptake with suction \citep{Feddes1978}, and the same idea underlies modern analyses of stomatal response to soil drying \citep{Carminati2020}. Yet no expression links that limit to $\alpha$, $n$ and $\ell$ in a form that an irrigation table could use.

The second question concerns the water once applied. In a series of numerical studies, \citet{Elmaloglou2007,Elmaloglou2008,Elmaloglou2009a,Elmaloglou2009b,Elmaloglou2010} showed that pulsing, discharge, spacing, hysteresis and the hydraulic properties of the soil all change the wetting pattern and the losses below the root zone under drip; \citet{Diamantopoulos2012VZJ} reviewed how flow departs from equilibrium; \citet{Weber2019} and \citet{Iden2021} extended the hydraulic functions to the dry range where films and vapour carry the flux; and \citet{Kumar2022} showed that site-specific schedules need site-specific hydraulic properties. Across that body of work, the share of an applied millimetre that leaves as transpiration, as evaporation or as deep percolation has been answered case by case, for one soil, one method and one dose at a time. How the shares depend jointly on the soil, the dose, the method and the season, and whether any of them follows a simple law, has not been mapped.

This study answers both questions with one theory and one data set: the dry-end asymptote of the Mualem conductivity, and an ensemble of 10\,190 Daisy simulations \citep{Hansen2012,Holbak2026} whose every output is openly archived (\cref{tab:data4}). Its objectives were (i) to derive a closed form for the reach $\hc$ of a suction trigger from $\alpha$, $n$ and $\ell$ with one constant, and to test it out of sample against interval-censored reaches of ten soils; (ii) to test whether a reach-scaled trigger $\htr/\hc$ places the irrigation responses of all soils on one axis; (iii) to quantify the marginal fate of an irrigation millimetre (transpiration, soil evaporation, other evaporation, drainage and storage) across eight soils, ten doses, ten seasons and three methods; and (iv) to identify the dependence of the drained share on the dose and translate the results into rules for sensor-scheduled irrigation.

%%PULL p4.tex fig:concept

\section{Materials and methods}
\subsection{The van~Genuchten--Mualem functions and the reach of a trigger}
With $h$ the suction (cm), the effective saturation and the relative conductivity are \citep{vanGenuchten1980,Mualem1976}
\begin{align}
S_e(h) &= \bigl[1+(\alpha h)^{n}\bigr]^{-m},\qquad m=1-1/n, \label{eq:vg}\\
K_r(h) &= S_e^{\,\ell}\Bigl[1-\bigl(1-S_e^{1/m}\bigr)^{m}\Bigr]^{2}, \label{eq:kr}
\end{align}
with $\ell=0.5$ as usual. For $\alpha h\gg1$, $S_e\simeq(\alpha h)^{-(n-1)}$ and $1-(1-S_e^{1/m})^{m}\simeq m\,S_e^{1/m}$, so
\begin{equation}
K_r(h)\;\simeq\;m^{2}\,(\alpha h)^{-[\,2n+\ell(n-1)\,]}, \qquad \alpha h\gg1.
\label{eq:asym}
\end{equation}
The exponent $2n+\ell(n-1)$ is 6.2 for the sand of this study and 2.6 for the silty clay loam: the conductivity of a coarse soil collapses over a narrow range of suction, that of a fine soil over many decades.

Roots take water from a layer at a rate limited by the conductance of the soil around them, which scales with $K$ \citep{Gardner1960,Carminati2020}; Daisy computes uptake from a soil--root potential difference over such a conductance \citep{Hansen1990,Abrahamsen2000}. As a layer dries, there is a suction beyond which the soil delivers too little water for the crop to dry it further within a season. If that point is reached at a fixed relative conductivity $\kappa$, \cref{eq:asym} gives the reach of a trigger at the sensor:
\begin{equation}
\hc \;=\; \frac{1}{\alpha}\left(\frac{m^{2}}{\kappa}\right)^{1/[\,2n+\ell(n-1)\,]}.
\label{eq:hc}
\end{equation}
A trigger $\htr<\hc$ can fire; a trigger $\htr>\hc$ is never reached and the crop is not irrigated, however dry it becomes. \Cref{eq:hc} has one constant, $\kappa$, shared by all soils, and it contains no saturated conductivity. Two alternatives that do (an absolute conductivity $K(\hc)=q_c$) and two that use only the retention curve (a fixed $S_e(\hc)$, or a power law in $n-1$) were fitted in the same way for comparison. The reach-scaled trigger is $x=\htr/\hc$.

\subsection{The marginal fate of irrigation}
For a run with irrigation $I$ and its rainfed twin on the same soil and weather (experiment E0), the change of each annual flux $X$ per unit of irrigation is
\begin{equation}
\eta_X=\frac{X-X_{0}}{I},\qquad X\in\{T,E_s,E_o,D\},\qquad \eta_T+\eta_{E_s}+\eta_{E_o}+\eta_D+\eta_S=1,
\label{eq:eta}
\end{equation}
where $T$ is transpiration, $E_s$ soil evaporation, $E_o=\ET-T-E_s$ the other evaporation (intercepted water, ponded water, litter and snow; these close $\ET$ within 1.1\,mm in the archived data), $D$ drainage at 1\,m and $\eta_S$ the change of storage and runoff, which closes the balance. Pooled over a set of runs, $\eta_X=\sum(X-X_0)/\sum I$. The share that drains was fitted, per soil, as
\begin{equation}
\eta_D(d)=\eta_\infty\bigl(1-e^{-d/d_D}\bigr),
\label{eq:dose}
\end{equation}
where $d$ is the dose per event, $\eta_\infty$ the plateau and $d_D$ the dose scale.

\subsection{Simulation ensemble and data}
Daisy solves the one-dimensional Richards equation \citep{Richards1931} with the functions of \cref{eq:vg,eq:kr}, coupled to a crop, a surface residue layer, an energy-based potential evapotranspiration and nitrogen turnover \citep{Hansen1991,Hansen2012,Holbak2026}. The ensemble has 17 experiments (E0 to E16), ten soils from sand to silty clay loam, 20 years of weather from Taastrup, Denmark (1980--1999) and a cereal rotation; 10\,041 of 10\,190 runs are valid and 149 are solver failures, which were excluded (\cref{tab:data4}). The irrigation scheduler of Daisy applies a dose whenever the suction at a sensor depth exceeds a trigger. The soils are described by their effective van~Genuchten parameters (\cref{tab:reachall}) and, for eight of them, by literature saturated conductivities \citep{Carsel1988}. Every output of every run, year by year, is in the repository folder \cslink{}, and the complete raw archive is on Zenodo \citep{HaliluData}.

%%PULL data_table_p4.tex tab:data4

Five experiments were used. \emph{E2} (197 runs) sets the trigger at 20 values from 50 to 10\,000\,cm on each soil (sensor at 20\,cm, surface irrigation). \emph{E3} (200 runs) moves the sensor from 5 to 100\,cm at a 300\,cm trigger. \emph{E10} (954 runs) crosses the dose (1 to 30\,mm per event, each applied in 12\,h) with ten irrigation windows on each soil, at a 300\,cm trigger and a 20\,cm sensor. \emph{E4} (96 runs) compares surface, overhead and subsurface application on each soil, and \emph{E15} (1000 runs, coarse loam) crosses the method with the rate, the dose and the maximum rooting depth. \emph{E0} gives the rainfed twin of every soil.

\subsection{Fitting and validation of the reach}
A trigger was said to fire on a soil if it brought at least 5\,mm of irrigation a year. For each soil the reach is then known as an interval: above the largest firing trigger and below the smallest non-firing one, or only bounded from below when even 10\,000\,cm fires (five soils). The closed forms were fitted by minimising the squared distance, in $\log_{10}$, between the prediction and the interval (zero inside the interval), and validated by leaving each soil out in turn (LOSO). The constant $\kappa$ of \cref{eq:hc} is therefore not measured but determined from experiment E2: one value is sought that places the predicted reach of all ten soils inside, or as close as possible to, their observed intervals; the minimum lies at $\kappa=3.47\times10^{-10}$. A bootstrap over soils gives a 95\,\% interval of $1.0\times10^{-10}$ to $2.8\times10^{-9}$, leaving one soil out gives $2.3$ to $7.3\times10^{-10}$, and a firing threshold of 1 to 40\,mm\,yr$^{-1}$ instead of 5 gives $2.7\times10^{-10}$ to $2.5\times10^{-9}$. Because $\hc\propto\kappa^{-1/E}$ with $E=2n+\ell(n-1)$, a tenfold change of $\kappa$ moves the reach by a factor of only 1.45 in sand and 2.45 in silty clay loam. The collapse of the irrigation responses on the reach-scaled axis was measured as the mean within-bin variance of the normalised irrigation (irrigation over its soil maximum) in ten quantile bins of $\log_{10}\htr$ or of $\log_{10}x$, relative to the total variance.

All analyses are in one script, \texttt{Manuscripts/reproducibility/analysis/p4\_reach.py}, which reads the data folder and writes every number of this paper to \texttt{Manuscripts/results/p4\_numbers.json} and every table and figure as a CSV file.

\section{Results}
\subsection{A suction trigger has a reach}
On every soil the irrigation fell as the trigger moved to drier values, but the fall ended in a cliff whose position differed by three orders of magnitude between soils (\cref{fig:trigger}a). In sand no trigger drier than 100\,cm fired; in loamy sand none drier than 200\,cm; in sandy loam none drier than 800\,cm; in loam and sandy clay loam none drier than 5\,000\,cm; the five other soils still irrigated at 10\,000\,cm. A 300\,cm trigger was therefore blind on sand and loamy sand: the crop stayed at its rainfed yield of 1.13 and 1.38\,Mg\,ha$^{-1}$, against 3.04 and 4.54\,Mg\,ha$^{-1}$ reached with a 50\,cm trigger (\cref{fig:trigger}d). Moving the sensor from 5 to 100\,cm did not cure the blindness: at a 300\,cm trigger, sand received no irrigation at any sensor depth and loamy sand at most 1.3\,mm a year (E3).

%%PULL p4.tex fig:curves

\subsection{The closed form against alternative rules}
\Cref{eq:hc} placed these cliffs with one constant (\cref{fig:curves,fig:trigger}c). With $\kappa=3.5\times10^{-10}$ the predicted reach was 199\,cm for sand, 426\,cm for loamy sand, 1\,618\,cm for sandy loam, 9\,209\,cm for loam and 7\,580\,cm for sandy clay loam, and above 5\,000\,cm for the five soils that still fired at 10\,000\,cm. Leaving each soil out, the prediction fell inside the observed interval for five soils and was never off by more than a factor of 2.3 (root-mean-square error 0.18 in $\log_{10}$; \cref{tab:laws}). The two-parameter power law in $n-1$ was slightly closer (0.14) but has no derivation, and a single effective saturation at the reach did worse (0.23, a factor of 3.2). A threshold on the absolute conductivity $K(\hc)$ was not supported: with the literature $K_s$ of eight soils the conductivity at the bracketing triggers spanned $10^{-4.8}$ to $10^{-8.4}$\,cm\,d$^{-1}$, more than three orders of magnitude, so the reach is set by the shape of the conductivity curve rather than by its level.

%%PULL p4.tex tab:laws

Whether the constant is worth having was judged against what is done without it (\cref{tab:base}). The common practice is to allow the sensor to trigger anywhere down to the wilting point of 15\,000\,cm. At that suction $K_r$ is $8\times10^{-22}$ in the sand and $9\times10^{-8}$ in the silty clay loam, so the rule treats a soil that cannot deliver any water as if it could, and it was wrong by a factor of 100 on the sand (RMSE 0.93). The best single suction for all ten soils (3\,140\,cm), the best proportion of $1/\alpha$ alone (the scale of the retention curve without its shape) and \cref{eq:hc} followed, in that order of improvement; each was fitted on all ten soils, which favours the baselines, and \cref{eq:hc} is also given with each soil left out. The baselines fell inside the interval only where the interval is open at the top (the five soils that still fire at 10\,000\,cm), which is why the count of soils inside is not used alone. As a consistency check (not an independent test, because the same intervals were used to fit $\kappa$), the last firing trigger of the five soils with a bracket lay at 0.47 to 0.66 of the predicted reach (\cref{fig:trigger}b), as expected if the reach is an upper limit that a trigger approaches without reaching.

%%PULL p4.tex tab:base

How $\kappa$ is found is shown in \cref{fig:kappa}: every trial value moves the predicted reach of all ten soils together (\cref{fig:kappa}a), and the summed squared distance in $\log_{10}$ to the observed intervals has its minimum, 0.196\,decades$^2$, at $\kappa=3.47\times10^{-10}$ (\cref{fig:kappa}b). The minimum of the fitting error is wide. With $\kappa=10^{-11}$, $10^{-10}$, $3.5\times10^{-10}$, $10^{-9}$ and $10^{-8}$ the root-mean-square error of the ten reaches was 0.32, 0.17, 0.14, 0.16 and 0.31 in $\log_{10}$ (fitted on all ten soils): a whole decade of $\kappa$ ($10^{-10}$ to $10^{-9}$) is within a factor of 1.25 of the optimum, which is why the bootstrap interval, the leave-one-soil-out range and the range for other firing thresholds all enclose it. A choice of $10^{-10}$ or $10^{-9}$ changes no conclusion of this study.

%%PULL p4.tex fig:trigger

\begin{figure}[t]
\centering
\includegraphics[width=\textwidth]{\figdir/paper4/fig6_kappa.pdf}
\caption{Determination of the constant $\kappa$ of \cref{eq:hc} from experiment E2. (a)~Predicted reach of each soil against the trial value of $\kappa$; at the right, the reach observed in Daisy (bars: interval between the last firing and the first non-firing trigger; arrows: only a lower bound is known, five soils); dashed line: the fitted value, $\log_{10}\kappa=-9.46$. (b)~The misfit, the squared distance in $\log_{10}$ between predicted reach and observed interval summed over the soils (zero inside the interval), with its minimum 0.196 at $\kappa=3.47\times10^{-10}$; dashed: the misfit with sand, loamy sand or sandy loam left out; grey: bootstrap 95\,\% interval over soils. Data: \texttt{p4\_fig6a\_reach\_vs\_kappa.csv}, \texttt{p4\_fig6b\_misfit.csv}.}
\label{fig:kappa}
\end{figure}

\subsection{One axis for all soils}
On the reach-scaled axis the end of irrigation fell between $x=0.5$ and 1 for every soil (\cref{fig:trigger}b). The within-bin variance of the normalised irrigation dropped from 54\,\% of the total on the raw trigger to 32\,\% on $x$ (Spearman $\rho=-0.83$ between $\log x$ and the normalised irrigation). What is left is the slope of the approach, which is steep in coarse soils and shallow in fine ones: the reach sets where irrigation ends, not how it fades. A practical rule follows: a trigger should lie below about $\hc/2$, and the 300\,cm of the tables satisfies this on eight of the ten soils but not on sand and loamy sand.

\subsection{Where an irrigation millimetre goes}
Pooled over the 744 irrigated runs of E10, a millimetre of irrigation became 0.38\,mm of transpiration, 0.25\,mm of soil evaporation, 0.19\,mm of other evaporation, 0.16\,mm of drainage below 1\,m and 0.03\,mm of storage and runoff (\cref{fig:fate}a). The transpired share was largest in sandy loam (0.64) and smallest in coarse sand (0.23), where 0.32 drained. In the five fine-textured soils the shares were close: 0.34 to 0.46 transpired, 0.15 drained, the rest evaporated.

The method changed the shares more than the soil did (\cref{fig:fate}b). Under overhead irrigation the extra transpiration was nil or negative ($-0.08$ in E4, $-0.01$ in E15): the water intercepted by the canopy evaporated (0.79 and 0.56 of each millimetre) and lowered the transpiration it replaced. Subsurface application transpired most (0.51 in E4, 0.40 in E15) but also drained most (0.23 and 0.28), because it bypasses the surface and wets the profile from within. Surface application lay between. In E15 the extra drainage under subsurface application rose to 92\,mm a year when the roots were limited to 30\,cm, against 56 to 60\,mm for any rooting depth of 45\,cm or more.

%%PULL p4.tex fig:fate

\subsection{The dose and the season}
The drained share rose with the dose and levelled off (\cref{fig:dose}a). \Cref{eq:dose} fitted every soil ($R^2$ 0.92 to 0.96 for seven soils, 0.66 for sandy loam) with a dose scale $d_D$ of 3.2 to 5.2\,mm, mean $4.1\pm0.6$\,mm, which was not ordered by any property of the soil (Spearman $\rho=0.07$ with $n$, $-0.24$ with the available water, $0.17$ with the water released between 100 and 300\,cm in the top 30\,cm). The plateau was soil-specific: 0.10 to 0.20 in seven soils and 0.45 in coarse sand. A dose of 1\,mm drained almost nothing (0.01 to 0.09), but 60 to 72\,\% of it evaporated, as expected for water held in the evaporating surface layer \citep{Lehmann2008,Shokri2011,Or2013} (\cref{fig:dose}c); the transpired share rose with the dose up to 15 to 20\,mm in most soils (\cref{fig:dose}b). The dose that maximised transpired minus drained water was 2\,mm on six of eight soils.

The irrigation window shifted the fate as strongly as the dose. Irrigating only from April to June transpired 0.64 of each millimetre and drained 0.07; irrigating the whole season (April to October) transpired 0.34 and drained 0.21; irrigating only from July to mid-September transpired 0.23. Early water was used by a growing canopy, whose water need rises with ground cover \citep{Allen1998,Doorenbos1979}; late water arrived when the cereal was ripening or harvested and was lost to evaporation and drainage.

%%PULL p4.tex fig:dose

\section{Discussion}
\subsection{Why one relative conductivity}
The constant $\kappa$ is the least obvious element of \cref{eq:hc}. A conductivity is the natural quantity because water reaches a root only through the soil around it, at a rate proportional to $K$; because $K_r$ falls by many orders of magnitude over a narrow range of suction (\cref{eq:asym}, \cref{fig:curves}b), any threshold on the conductivity becomes a sharp threshold on the suction, which is the cliff seen in experiment E2. A threshold on the water content is blunt, because $S_e$ changes by a few tenths over the same range, and this is why the fixed-saturation law was the worst of the three (\cref{tab:laws}). The threshold is relative and not absolute because the saturated conductivity of a soil is not fixed by its retention curve and differs widely between sources for the same texture, whereas $K_r$ is dimensionless and fixed by $\alpha$, $n$ and $\ell$; the absolute form did not hold. The relative form also scales correctly by construction: multiplying $\alpha$ by a factor shifts the retention and conductivity curves along the suction axis by the inverse factor, and \cref{eq:hc} moves the reach by exactly that factor. Finally, one constant serves all soils because the soils differ in the shape of the curves, which enters through $n$ and $\ell$, whereas $\kappa$ stands for what did not differ between them: the crop, the climate and the sensor (root density at the sensor, transpiration demand, lowest root water potential). For another crop or climate it is redetermined in the same way from a trigger experiment on a few soils.

\subsection{What the closed form adds to the van~Genuchten curve}
The van~Genuchten--Mualem pair describes a soil at every suction; \cref{eq:hc} says which suctions a cropped soil actually visits at the depth of a sensor. Its exponent, $2n+\ell(n-1)$, is the slope of the conductivity on a log--log scale at the dry end, so the reach is short where the conductivity collapses fast (large $n$) and long where it fades slowly (small $n$), and $\alpha$ only rescales it. This complements the wet-end criteria of field capacity derived from the same parameters \citep{Twarakavi2009,Assouline2014}: together they bracket the working range of a scheduled soil. That a relative and not an absolute conductivity sets the reach suggests that the limit is reached when the conductivity has fallen a fixed number of decades below its saturated value, which a field test can check.

\subsection{The fate of a millimetre and the pulsing question}
The pulsed and intermittent drip simulations of \citet{Elmaloglou2007,Elmaloglou2008,Elmaloglou2009b} found that the application regime changes the losses below the root zone and that the effect depends on the soil. Across ten soils, ten doses, ten windows and three methods, the present data say how: the dose sets the drained share on a scale of about one day of crop water use, the soil sets the plateau, the method sets how much of the water can be transpired at all, and the season sets how much the crop can use. A dose smaller than one day of evapotranspiration is used before it can move below the roots, whatever the soil, while the share of a large dose that is lost depends on how fast the soil drains once wetted, which the hydraulic parameters set. Small doses do suppress drainage, but in a cropped field with an evaporating surface the water they save is lost to evaporation instead; the best compromise in these runs was a dose of about 2\,mm. This is a marginal accounting and does not replace the event-scale wetting patterns of the earlier work, which a one-dimensional analysis of annual means cannot resolve.

\subsection{Practical rules for sensor-scheduled irrigation}
Three rules follow. First, compute $\hc$ from \cref{eq:hc} for the soil and keep the trigger below $\hc/2$; on sands this means triggers of 50 to 100\,cm, not the 300 to 600\,cm of tables. Second, choose the dose against the daily crop water use: doses well below it evaporate, doses well above it drain at the soil's plateau. Third, judge overhead irrigation by its effect on transpiration, which can be nil, and not by the water delivered. These rules complement the established practice of deficit and scheduled irrigation \citep{Fereres2007,Geerts2009,Pereira2012} and the FAO-56 root-zone balance \citep{Allen1998}. The reach of every soil and the trigger law are implemented in the open decision-support software SWAFL Desktop 1.1.0, which draws the trigger of every soil against its reach and runs new trigger, dose and method experiments with Daisy on soils of the user's choice \citep{HaliluManual}; screenshots are given in the Supplementary Material.

\subsection{Limitations and the experiments that would test the claims}
The results are from one process model, one weather series, one cereal rotation, one-dimensional columns, equilibrium Richards flow without hysteresis or water repellency \citep{Elmaloglou2009b,Diamantopoulos2013AWR,Diamantopoulos2013JH}, and effective van~Genuchten parameters; the Mualem connectivity was not varied, the sensor was at 20\,cm in E2 and E10, and the firing threshold of 5\,mm a year is a choice (the cliffs are steep, so the intervals change little with it). Five of ten reaches are only bounded from below. The constant $\kappa$ is a property of the model's uptake and should be measured, not assumed, in the field. The fates are 20-year means and do not resolve events.

Under fast wetting and drying the water content lags the equilibrium of the retention curve, $\partial\theta/\partial t=[\theta_{\mathrm{eq}}(h)-\theta]/\tau$ \citep{RossSmettem2000,Hassanizadeh2002,Diamantopoulos2012VZJ,Diamantopoulos2012WRR}, and a sensor then reads a suction that is not the one of \cref{eq:vg}. For a schedule with interval $t_i$ between applications, the dimensionless lag $\Pi_\tau=\tau/t_i$ says when this matters: for daily or multi-day scheduling and $\tau$ of minutes to hours, $\Pi_\tau\ll1$ and the equilibrium reach applies; for pulses of minutes to hours, as in pulsed drip, $\Pi_\tau\sim1$ and the sensor reading, the effective $\kappa$ and the drained share become rate-dependent. Daisy solves the equilibrium Richards equation, so this criterion is stated for later testing. A field test would place matric-potential sensors on two coarse and two fine soils, schedule plots at triggers on both sides of $\hc$ and at doses of 1, 4 and 15\,mm, and measure drainage with lysimeters \citep{Groh2020}.

\section{Conclusions}
A suction trigger has a reach, the largest suction a crop can impose at the sensor, and it follows from the van~Genuchten--Mualem parameters by one closed form with one constant, $\hc=\alpha^{-1}(m^{2}/\kappa)^{1/[2n+\ell(n-1)]}$ with $\kappa=3.5\times10^{-10}$, valid out of sample within a factor of 2.3 over ten soils from sand to silty clay loam. Beyond the reach a trigger never fires, and the common 300\,cm trigger left sand and loamy sand rainfed. Of each irrigation millimetre, 38\,\% was transpired, 44\,\% evaporated and 16\,\% drained; the drained share rose with the dose on a soil-independent scale of about 4\,mm towards a soil-specific plateau, and overhead irrigation transpired nothing extra. Triggers should therefore be set against the reach of the soil, and doses against the daily water use of the crop.

\section*{CRediT authorship contribution statement}
\textbf{Naziru Halilu:} Conceptualization, Methodology, Software, Formal analysis, Investigation, Data curation, Visualization, Writing -- original draft, Writing -- review \& editing.

\section*{Declaration of competing interest}
The author declares that he has no known competing financial interests or personal relationships that could have appeared to influence the work reported in this paper.

\section*{Funding}
This research did not receive any specific grant from funding agencies in the public, commercial, or not-for-profit sectors.

\section*{Data availability}
The complete Daisy ensemble outputs (all 10\,190 runs) are archived on Zenodo (\doihl{10.5281/zenodo.23043626}). Every log of every run, year by year, and the one-row-per-run summary are in the repository folder \cslink{} (\url{\csurl}). The analysis script (\texttt{Manuscripts/reproducibility/analysis/p4\_reach.py}) and the software SWAFL Desktop 1.1.0 (MIT licence) are in the project repository (\url{\repourl}).

\section*{Acknowledgements}
The Daisy model is developed and maintained at the University of Copenhagen. The author thanks the maintainers of the open-source libraries used.

\appendix
\section{Reach of every soil}\label{app:reach}
%%PULL p4.tex tab:reachall
The two soils whose observed lower bound exceeds the prediction (coarse sand and coarse loam) are the two without a literature saturated conductivity in the set used here; their reach is underpredicted by at least a factor of 1.9 and 1.3, which a second constant (the power law of \cref{tab:laws}) only partly removes.

\section*{Supplementary material}
Supplementary material associated with this article is available online. \textbf{Supplementary Material S1} gives the complete text of the study in the order of this article, with every section, figure, table, equation and appendix (numbered S1, S2, \ldots), including the root-zone depletion view of the trigger law, the flux diagrams, the controllability of each outcome and the screenshots of the software. \textbf{Supplementary Material S2} gives the derivations, the data and the mechanisms behind every table and figure, and the map from every number to its script and data file.

\bibliographystyle{elsarticle-harv}
\bibliography{../tex/refs}
\end{document}
